\section*{第 \ref{ch:b3:orogenic-wedges} 章——逆冲带与造山楔}

\begin{solution}{exo:b3:orogenic-wedges:1}
$\sin\ang{30} = 0.5$：$K = 2\times0.5/0.5 = 2$，$\sigma'_1/\sigma'_3 = 1 +
K = 3$。
\end{solution}

\begin{solution}{exo:b3:orogenic-wedges:2}
$\phi = \arctan0.6 = \ang{31.0}$，$\sin\phi = 0.514$，$K = 2.12$。
在水平底面上，$\alpha = 0.3/3.12 = \qty{0.0962}{rad} = \ang{5.5}$。
在倾角为 $\ang{3} = \qty{0.0524}{rad}$ 的底面上，$\alpha + \beta = 0.352/3.12 =
\qty{0.113}{rad} = \ang{6.5}$，因此 $\alpha = \ang{3.5}$。
\end{solution}

\begin{solution}{exo:b3:orogenic-wedges:3}
深度 $= (T - 10)/30$：磷灰石（U--Th）/He 为 1.3 至 \qty{2.3}{km}；
磷灰石裂变径迹为 2.7 至 \qty{3.7}{km}；锆石裂变径迹 $270/30 = \qty{9.0}{km}$。
在 \qty{5}{km} 深处，岩石温度为 $10 + 150 = \qty{160}{\celsius}$：锆石
（U--Th）/He 的封闭温度为 160 至 \qty{200}{\celsius}，对应深度为 5.0 至
\qty{6.3}{km}。
\end{solution}

\begin{solution}{exo:b3:orogenic-wedges:4}
水平长度 $2/\tan\ang{30} = \qty{3.46}{km}$；
沿倾向长度 $2/\sin\ang{30} = \qty{4.0}{km}$。岩层被\omterm{def:b3:orogenic-wedges:thrust-geometry}{断坡}高度抬升
\qty{2}{km}。截断点须穿过整个\omterm{def:b3:orogenic-wedges:thrust-geometry}{断坡}，滑距为
\qty{4.0}{km}。
\end{solution}

\begin{solution}{exo:b3:orogenic-wedges:5}
$\phi = \arctan1.03 = \ang{45.8}$，$K = 5.08$；$\beta = \qty{0.105}{rad}$。
当 $\lambda = 0.4$ 时：$\alpha + \beta = (0.105 + 0.85\times0.6)/(1 +
5.08\times0.6) = 0.615/4.05 = \qty{0.152}{rad} = \ang{8.7}$，$\alpha =
\ang{2.7}$。当 $\lambda = 0.9$ 时：$0.190/1.51 = \qty{0.126}{rad} = \ang{7.2}$，
$\alpha = \ang{1.2}$。提高流体压力会降低\omterm{def:b3:orogenic-wedges:critical-taper}{临界楔角}：
楔体变为超临界状态，沿其底面向前滑动而不增厚，
其前缘向前推进，直至增生与侵蚀使其回到
新的楔角。
\end{solution}

\begin{solution}{exo:b3:orogenic-wedges:6}
$S = A/r = 6.0/1.5 = \qty{4.0}{km}$。若按\omterm{def:b3:orogenic-wedges:thrust-folds}{滑脱褶皱}解释，则 $S = A/h =
6.0/3.0 = \qty{2.0}{km}$：仅为真实值的一半。两种解释须
通过褶皱形态及是否存在\omterm{def:b3:orogenic-wedges:thrust-geometry}{断坡}加以区分。
\end{solution}

\begin{solution}{exo:b3:orogenic-wedges:7}
$K = 2.12$，$1 + K(1 - 0.4) = 2.27$。方程给出 $\mu_b(1 -
\lambda_b) = (\alpha + \beta)\times2.27 - \beta$，其中 $\beta =
\qty{0.0349}{rad}$。侏罗山：$0.0611\times2.27 - 0.0349 = 0.10$。南侧：
$0.218\times2.27 - 0.0349 = 0.46$。静水压力下的拜尔利值：
$0.85\times0.6 = 0.51$。南侧底面接近拜尔利值；
侏罗山底面的强度几乎只有南侧的五分之一，符合岩盐中滑脱面的预期：
岩盐摩擦力低，且流体并未超压。
\end{solution}

\begin{solution}{exo:b3:orogenic-wedges:8}
$\kappa = \qty{e-6}{m^2/s} = \qty{31.6}{km^2/Ma}$，且 $v =
\qty{1}{km/Ma}$：$\mathrm{Pe} = 30/31.6 = 0.951$。因子 $0.951/(1 -
e^{-0.951}) = 1.55$ 使梯度由 30 提高到 \qty{46.5}{\celsius/km}。
对于 \qty{65}{\celsius}：$1 - e^{-\mathrm{Pe}\,z/L} = (55/900)(1 - e^{-0.951})
= 0.0375$，故 $\mathrm{Pe}\,z/L = 0.0382$，$z = \qty{1.21}{km}$，而不计侵蚀时为
$55/30 = \qty{1.83}{km}$。
\end{solution}

\begin{solution}{exo:b3:orogenic-wedges:9}
$\mathrm dh/\mathrm dt = \qty{0.8}{km}/\qty{1.2}{Ma} = \qty{0.67}{mm/\text{年}}$。
取 $f = 0.5$，$v = 0.67\times(1 - 0.5) = \qty{0.33}{mm/\text{年}}$。
\end{solution}

\begin{solution}{exo:b3:orogenic-wedges:10}
取 $\alpha = 0$，$u = 1 - \lambda$：$\beta(1 + Ku) = \beta + \mu_b u$，
故 $\beta Ku = \mu_b u$，$\beta^* = \mu_b/K$，与 $u$ 无关。对台湾而言，
$0.85/5.08 = \qty{0.167}{rad} = \ang{9.6}$。在更陡的底面上，临界表面坡度为负值：
任何表面水平或朝前陆倾斜的楔体都已处于超临界状态，
会沿其底面滑动而不增厚。逆冲带无法
在此发育其正常形态。
\end{solution}

\begin{solution}{exo:b3:orogenic-wedges:11}
当 $\Lambda \to \infty$ 且 $v \to 0$ 时，$m \to 0$，故 $e^{-mz_c} \to 1$，
地温线为直线，$G_0 = G(z_c) = G$：$f \to (G - \Gamma)/G = 1 -
\Gamma/G = 1 - 6.5/30 = 0.78$。此时坡度为 $v/(1 - f) = vG/\Gamma =
4.6v$。对于 $v = \qty{1}{mm/\text{年}}$，$v/2\kappa = \qty{0.0158}{km^{-1}}$。在
\qty{5}{km}处：$m = 0.0158 + \sqrt{0.0158^2 + 1.257^2} = \qty{1.273}{km^{-1}}$，
$e^{-1.273\times1.21} = 0.215$，$f = 40.0\times0.215/44.7 = 0.19$。在
\qty{20}{km}处：$m = 0.0158 + 0.3146 = \qty{0.330}{km^{-1}}$，$e^{-0.399} =
0.671$，$f = 40.0\times0.671/44.7 = 0.60$。
\end{solution}

\begin{solution}{exo:b3:orogenic-wedges:12}
该剖面为底为 $W$、高为 $W\tan(\alpha + \beta)$ 的三角形：
$A = W^2\tan(\alpha + \beta)/2$。由 $\mathrm dA/\mathrm dt = F_a$ 且 $A =
0$ 于 $t = 0$ 时取得，得 $A = F_at$，由此求出 $W$。取 $\tan\ang{6} = 0.105$：
$W = \sqrt{2\times40\times2/0.105} = \qty{39}{km}$（\qty{2}{Ma} 后）；要达到 \qty{60}{km}，$t =
60^2\times0.105/(2\times40) = \qty{4.7}{Ma}$。
\end{solution}

\begin{solution}{pb:b3:orogenic-wedges:1}
\textbf{1.} $4.5 - 1.5 = \qty{3.0}{km}$ 高；$3.0/\tan\ang{25} =
\qty{6.43}{km}$ 水平长；$3.0/\sin\ang{25} = \qty{7.10}{km}$ 沿
倾向长。

\textbf{2.} 由\omterm{def:b3:orogenic-wedges:thrust-geometry}{断坡}高度可得，为 \qty{3.0}{km}。

\textbf{3.} $S = A/r = 16.5/3.0 = \qty{5.5}{km}$。

\textbf{4.} 面积同为 \qty{16.5}{km^2}：在\omterm{def:b3:orogenic-wedges:thrust-folds}{断层转折褶皱}中，
每一岩层位于上部\omterm{def:b3:orogenic-wedges:thrust-geometry}{断坪}以上的多余面积为 $Sr$。若按\omterm{def:b3:orogenic-wedges:thrust-folds}{滑脱褶皱}计算，$16.5/4.0 =
\qty{4.1}{km}$：偏小四分之一。

\textbf{5.} $4\times(6.0 - 2.5) = \qty{14.0}{km}$，为前缘\omterm{def:b3:orogenic-wedges:thrust-geometry}{断坡}上
\qty{5.5}{km} 滑距的两倍半。

\textbf{6.} 前缘岩片，即最靠近前陆者：每条新\omterm{def:b3:orogenic-wedges:thrust-geometry}{断坡}
于前一条\omterm{def:b3:orogenic-wedges:thrust-geometry}{断坡}之前的下盘中形成。每个岩片沿其\omterm{def:b3:orogenic-wedges:thrust-geometry}{断坡}被抬升，
正如\omterm{def:b3:orogenic-wedges:thrust-folds}{断层转折褶皱}抬起其上盘一样，同时抬起顶板逆冲断层及其上的盖层，
使顶板在整个堆叠体上方发生褶皱。

\textbf{7.} $\alpha = \arctan(2.1/60) = \ang{2.00}$；$\beta =
\arctan(4.2/60) = \ang{4.00}$；楔角 \ang{6.0}。

\textbf{8.} $2.4 + 8.2 = \qty{10.6}{km}$，小于约 \qty{15}{km}：基底
仍具摩擦性，故适用库仑理论。

\textbf{9.} $\phi = \ang{45.8}$，$\sin\phi = 0.717$，$K = 5.08$。干燥情形：
$\alpha + \beta = (0.0699 + 0.85)/6.08 = \qty{0.151}{rad} = \ang{8.7}$，
$\alpha = \ang{4.7}$。

\textbf{10.} $\lambda = 1000/2500 = 0.40$：$(0.0699 + 0.51)/4.05 =
\qty{0.143}{rad} = \ang{8.2}$，$\alpha = \ang{4.2}$。观察到的 \ang{2.0}
会使楔体强烈亚临界，即增厚而非前进：
流体必定处于超压状态。

\textbf{11.} 令 $T = \qty{0.1049}{rad}$，$u = 1 - \lambda$，$T(1 + Ku) =
\beta + \mu_bu$，可得 $u = (\beta - T)/(TK - \mu_b) = (0.0699 -
0.1049)/(0.5327 - 0.85) = 0.110$：$\lambda = 0.89$。

\textbf{12.} $\sigma_v = 2500\times9.81\times10\,600 = \qty{260}{MPa}$；$P =
0.89\times260 = \qty{231}{MPa}$；静力的压强 $1000\times9.81\times10\,600 =
\qty{104}{MPa}$。流体承载的压强比同深度水柱多出
\qty{127}{MPa}。

\textbf{13.} $(0.0699 + 0.85\times0.3)/(1 + 5.08\times0.3) = 0.325/2.52 =
\qty{0.129}{rad} = \ang{7.4}$，$\alpha = \ang{3.4}$。楔体坡角为 \ang{2.0}，
已处于亚临界状态：它停止前进，并通过\omterm{def:b3:orogenic-wedges:thrust-geometry}{无序逆冲断层}在其前缘后方增厚，
直至其地表坡度达到 \ang{3.4}。

\textbf{14.} $F_a = \qty{10}{km/Ma}\times\qty{4.0}{km} = \qty{40}{km^2/Ma}$，
即每千米造山带 \qty{4.0e4}{m^3/\text{年}}。

\textbf{15.} $\bar e = 40/60 = \qty{0.67}{km/Ma} = \qty{0.67}{mm/\text{年}}$。

\textbf{16.} $40\times100 = \qty{4000}{km^3/Ma} = \qty{4.0e6}{m^3/\text{年}}$，或
$\qty{1.0e10}{kg/\text{年}}$（每年 1000 万吨）。

\textbf{17.} $\mathrm dW/\mathrm dt = F_a/(W\tan(\alpha + \beta)) =
40/(60\times0.105) = \qty{6.3}{km/Ma}$。

\textbf{18.} 否：稳态时，侵蚀仍须移走整个
输入通量，故平均速率仍为 \qty{0.67}{mm/\text{年}}。但在后缘之下加入的岩石
在该处被抬升：侵蚀与折返将最快发生
于后缘上方，那里应见到最年轻的冷却年龄。

\textbf{19.} $v = \qty{0.667}{km/Ma}$：$\mathrm{Pe} = 0.667\times30/31.6 =
0.634$；$G_0 = 30\times0.634/(1 - e^{-0.634}) = 30\times1.35 =
\qty{40.5}{\celsius/km}$。

\textbf{20.} $1 - e^{-\mathrm{Pe}\,z/L} = (55/900)(1 - e^{-0.634}) = 0.0287$，
故 $\mathrm{Pe}\,z/L = 0.0291$，且 $z_c = \qty{1.38}{km}$；$G(z_c) =
40.5\,e^{-0.0291} = \qty{39.3}{\celsius/km}$。

\textbf{21.} $1.38/0.667 = \qty{2.07}{Ma}$，即样品在
\qty{2.00}{km} 处的年龄：在平均高程处，等温线的挠曲为零。

\textbf{22.} $v/2\kappa = 0.667/63.1 = \qty{0.0106}{km^{-1}}$，$2\pi/\Lambda =
\qty{0.314}{km^{-1}}$：$m = 0.0106 + 0.3143 = \qty{0.325}{km^{-1}}$；
$e^{-0.325\times1.38} = 0.639$；$f = (40.5 - 6.5)\times0.639/39.3 = 0.552$。

\textbf{23.} 平均高程为 \qty{2.00}{km}，平均年龄
为 \qty{2.068}{Ma}；$\sum\Delta t\,\Delta h = 0.823$，且 $\sum\Delta t^2 =
0.552$，故斜率为 \qty{1.49}{km/Ma}：视速率为
\qty{1.49}{mm/\text{年}}。

\textbf{24.} $v = 1.49\times(1 - 0.552) = \qty{0.667}{mm/\text{年}}$。

\textbf{25.} \omterm{def:b3:orogenic-wedges:thermochronology}{热年代计}与通量平衡结果相符：楔体
处于稳态；若将视斜率照面值取用，
则会将速率高估一倍以上。凯斯特尔
楔体的折返速率为 \textbf{\qty[text-series-to-math=true]{0.67}{mm/\text{年}}}。
\end{solution}
